\chapter{Substance misuse}
\index{Substance misuse}

\medskip
\noindent\textit{Educational reference; always seek professional advice for individual care.}

\section*{Definition and Overview}
Substance misuse, sometimes called drug misuse or substance use disorder, describes the harmful or hazardous use of alcohol, prescription medicines, or illicit drugs. It covers a broad range of patterns, from occasional risky use to severe dependence, and includes legal substances used in harmful ways as well as illegal ones. The term is deliberately broad because the harms can arise from a single episode of intoxication, from regular heavy use, or from the physical and psychological dependence that develops over time. Substance misuse is not a moral failing or a sign of weak character; it is recognised as a health condition that affects the brain, behaviour, and body, and it responds to care and support in the same way other chronic conditions do.

The substances involved are many: alcohol, tobacco, cannabis, stimulants such as methamphetamine and cocaine, opioids such as heroin and prescription painkillers, sedatives and benzodiazepines, hallucinogens, and inhalants. Each has its own effects, risks, and withdrawal patterns, but they share common threads: impaired judgement, damage to physical and mental health, strain on relationships, and risk of dependence. Substance misuse is common in many countries and worldwide, and it affects people of all ages, backgrounds, and incomes. This chapter gives an overview of the condition, how to recognise it, how it affects the body and mind, and the treatment and support available. Throughout, the aim is to inform without judgement, because people who use substances are far more likely to reach out for help when they feel safe and respected.

\section*{Epidemiology}
Substance misuse is among the most common health problems in many countries. Alcohol is the most widely used and misused drug, with a substantial proportion of adults drinking at levels that exceed the national guidelines, and a smaller but significant minority drinking heavily. Tobacco remains a major cause of preventable illness and death even though smoking rates have fallen over decades. Cannabis is the most commonly used illicit drug, and a meaningful minority of people use stimulants such as methamphetamine, mostly occasionally but some heavily. Opioid misuse, involving both heroin and pharmaceutical opioids, is a serious concern, and in recent years the harms from prescribed and over-the-counter medicines, including benzodiazepines, have also become clearer.

The burden falls unevenly. Young adults, people living with disadvantage, and people with mental illness or trauma histories may face higher rates of harm. Substance misuse commonly coexists with depression, anxiety, and trauma, and the two problems reinforce each other. The pattern has shifted over time: rates of smoking have declined in many places, while concerns have grown about stimulant harms, opioid dependence, and increasingly potent products. Substance misuse contributes to hospital admissions, accidents, violence, family breakdown, and suicide, but effective treatment is available.

\section*{Aetiology and Pathophysiology}
Substance misuse develops from the interaction of the drug, the person, and their environment. All addictive substances act on the brain's reward system, particularly on the chemical dopamine, producing feelings of pleasure and reinforcement that make the drug more likely to be repeated. Over time the brain adapts: it becomes less responsive to the drug, so more is needed to produce the same effect (tolerance), and it becomes dependent on the drug to function normally, so stopping causes withdrawal symptoms. This is a physical change in brain circuitry, which is why dependence is a medical condition rather than a simple choice.

Why some people become dependent and others do not is influenced by genetics, age of first use, mental health, trauma, and social circumstances. People who use substances to cope with pain, stress, or difficult emotions are especially vulnerable, as are young people whose brains are still developing. Environmental factors, including poverty, family substance use, availability of drugs, and exposure to advertising, shape the risk for whole communities.

\begin{itemize}
    \item \textbf{Alcohol:} the most widely misused substance; heavy use damages the liver, brain, heart, and pancreas, and causes dependence
    \item \textbf{Tobacco and nicotine:} powerfully addictive; the leading preventable cause of cancer, heart disease, and lung disease
    \item \textbf{Cannabis:} can impair memory, concentration, and motivation, and worsen anxiety and psychosis in vulnerable people
    \item \textbf{Stimulants:} methamphetamine, cocaine, and related drugs overstimulate the brain, causing agitation, paranoia, heart strain, and psychosis
    \item \textbf{Opioids:} heroin and pharmaceutical opioids relieve pain and produce euphoria but carry a high risk of dependence and fatal overdose
    \item \textbf{Benzodiazepines and sedatives:} calming medicines that cause dependence and dangerous withdrawal if stopped abruptly
    \item \textbf{Hallucinogens and inhalants:} unpredictable effects including hallucinations, accidents, and brain damage, particularly with inhalants
\end{itemize}

\section*{Clinical Features}
The features of substance misuse depend on the substance and the pattern of use, but several changes tend to appear across all drugs.

\begin{itemize}
    \item \textbf{Tolerance:} needing more of the substance to get the same effect
    \item \textbf{Withdrawal:} unpleasant physical and mental symptoms when the substance wears off, such as anxiety, sweating, tremors, nausea, or seizures
    \item \textbf{Loss of control:} using more, or for longer, than intended, and unsuccessful attempts to cut down
    \item \textbf{Craving:} a strong urge to use that dominates thinking
    \item \textbf{Neglect of responsibilities:} declining work, study, or family roles because of use
    \item \textbf{Continuing despite harm:} using on even when it causes health, financial, or relationship problems
    \item \textbf{Risky use:} driving while affected, sharing needles, using in dangerous settings, or mixing substances
    \item \textbf{Changes in mood and behaviour:} irritability, secrecy, mood swings, withdrawal from friends and family
    \item \textbf{Physical effects:} weight change, poor sleep, injection-site infections, accidents, and poor general health
\end{itemize}

\section*{Differential Diagnosis}
\begin{itemize}
    \item \textbf{Depression and anxiety:} low mood, poor sleep, and withdrawal from others can be a mental illness rather than substance misuse, or the two can coexist and need both to be treated
    \item \textbf{Bipolar disorder:} impulsive spending, risk-taking, and elevated mood can resemble stimulant intoxication, and stimulants can trigger mania
    \item \textbf{Psychotic illness:} paranoia and hallucinations occur in schizophrenia and in stimulant or cannabis psychosis; careful history-taking and observation distinguish them
    \item \textbf{ADHD and other neurodevelopmental conditions:} restlessness and risk-taking can be mistaken for substance effects, and untreated ADHD is itself a risk factor for substance misuse
    \item \textbf{Chronic pain and fatigue conditions:} people using prescribed opioids to cope may be mislabelled as misusing, or legitimate use can drift into dependence
    \item \textbf{Dementia and cognitive decline:} confusion, poor judgement, and memory problems can mimic intoxication in older people
    \item \textbf{Thyroid and other medical disorders:} weight change, tremor, and mood disturbance can be medical in origin
\end{itemize}

\section*{When to Seek Medical Care}
Anyone who is concerned about their own or someone else's substance use should speak to a general practitioner, who can assess the pattern of use, screen for complications, and arrange support. Many people find that simply talking about it with a trusted professional is a relief. Early help prevents the harms that come with long-term use.

If you or someone you care about is experiencing distressing thoughts or feelings, contact local emergency services in an immediate crisis, or seek help from a local mental-health, alcohol-and-drug, youth, or culturally safe support service.

\textbf{Contact your local emergency services for:}
\begin{itemize}
    \item Severe intoxication with difficulty breathing, collapse, seizures, or unresponsiveness
    \item Signs of overdose, such as slow or stopped breathing, pinpoint pupils, or blue lips
    \item Severe agitation, paranoia, or violence, particularly after stimulants
    \item Thoughts of self-harm or suicide, or a serious accident or injury while affected
\end{itemize}

\section*{Diagnosis and Investigations}
\begin{itemize}
    \item \textbf{Clinical history:} a respectful, confidential discussion of what is used, how much, how often, the impact on life, and any attempts to stop; screening questions such as the AUDIT for alcohol are commonly used
    \item \textbf{Mental state assessment:} mood, thinking, memory, and any symptoms of psychosis or suicidality are assessed
    \item \textbf{Physical examination:} a check for signs of intoxication or withdrawal, injection marks, and the physical complications of the particular substance
    \item \textbf{Blood tests:} liver function, kidney function, and blood counts detect organ damage; blood alcohol or drug levels may be measured
    \item \textbf{Drug testing:} urine or saliva tests for substances can confirm use but are usually used in treatment settings rather than as a barrier to care
    \item \textbf{Infectious disease screening:} hepatitis B and C and HIV testing are important for people who inject drugs
    \item \textbf{Specialist assessment:} addiction medicine specialists and alcohol and drug services provide comprehensive assessment and treatment planning
\end{itemize}

\section*{Management}
\begin{itemize}
    \item \textbf{Recognition and motivation:} non-judgemental support that helps the person recognise the impact of their use and builds motivation to change; brief GP advice is effective for many people
    \item \textbf{Talking therapies:} cognitive behaviour therapy and related approaches help identify triggers, develop coping strategies, and prevent relapse
    \item \textbf{Medication-assisted treatment:} medicines such as methadone and buprenorphine for opioid dependence reduce craving and overdose risk and allow stability while other life changes happen
    \item \textbf{Management of dependence and withdrawal:} medically supervised withdrawal (detoxification) makes stopping safer and more comfortable, and is often the first step in treatment
    \item \textbf{Treatment of coexisting conditions:} depression, anxiety, trauma, and chronic pain must be treated alongside the substance misuse for either to improve
    \item \textbf{Harm reduction:} practical measures that reduce immediate danger, including naloxone for opioid overdose, needle and syringe programs, and supervised injecting rooms
    \item \textbf{Counselling and support groups:} peer support and family involvement strengthen recovery and reduce isolation
    \item \textbf{Residential rehabilitation:} structured live-in programs help people who need a period away from their usual environment
    \item \textbf{Ongoing relapse prevention:} recovery is often a long process with setbacks; regular review and continuing support make success more likely
\end{itemize}

\section*{Complications}
\begin{itemize}
    \item \textbf{Overdose:} the most dangerous acute complication, especially with opioids, alcohol, and sedatives
    \item \textbf{Organ damage:} liver cirrhosis, pancreatitis, heart disease, lung disease, and kidney damage with long-term use
    \item \textbf{Mental health problems:} depression, anxiety, psychosis, and an increased risk of suicide
    \item \textbf{Infections:} hepatitis B and C, HIV, and bacterial infections of the heart and injection sites
    \item \textbf{Injury:} accidents, falls, drowning, and burns, particularly with intoxication
    \item \textbf{Pregnancy harms:} substance misuse in pregnancy affects the developing baby and can cause withdrawal after birth
    \item \textbf{Social harms:} relationship breakdown, financial loss, unemployment, and legal problems
    \item \textbf{Withdrawal complications:} seizures and delirium can occur when heavy alcohol or sedative use is stopped suddenly
\end{itemize}

\section*{Prognosis and Outcomes}
Recovery from substance misuse is realistic and common, but it is usually not a single event; it is a process that takes time and often involves setbacks. People who engage with treatment do substantially better than those who do not. The best outcomes occur when the person is supported to make their own decisions, when coexisting mental health and social problems are addressed, and when treatment continues for long enough, because relapse is commonest in the first months after stopping. Medication-assisted treatment for opioid dependence in particular is associated with greatly reduced overdose deaths and improved health, employment, and family outcomes. Even after many years of heavy use, the brain and body show remarkable capacity for recovery once the substance is stopped. The overarching message is hopeful: substance misuse is treatable, help is effective, and many people in recovery lead full, healthy, and productive lives. The earlier someone reaches out, the more of their life can be protected.

\section*{Prevention and Self-Care}
\begin{itemize}
    \item Drink within the local guidelines, which recommend no more than ten standard drinks a week and no more than four on any one day
    \item Get support to quit smoking; nicotine replacement and quit services double the chance of success
    \item Do not mix substances, and never drink or use drugs when driving, boating, or operating machinery
    \item If you inject drugs, always use a new needle from a needle and syringe program and never share equipment
    \item Talk openly with your doctor about your substance use; they cannot help with what they do not know about
    \item Seek help early from a GP or an alcohol and drug service; early intervention is far easier than treating severe dependence
    \item If you are close to someone who uses substances, learn about overdose response, including naloxone, and keep a support network for yourself
    \item Be kind to yourself; relapse is part of recovery for many people and is a reason to continue, not to give up
\end{itemize}

\section*{Sources and Further Reading}
The following organisations provide reliable information and support for substance misuse.

\begin{itemize}
    \item World Health Organization. \textit{Alcohol and drug information and support.} https://www.who.int/health-topics/drugs-psychoactive
    \item NHS. \textit{Drug addiction: getting help and treatment.} https://www.nhs.uk/
    \item Mayo Clinic. \textit{Drug addiction: symptoms and causes.} https://www.mayoclinic.org/
    \item National Institute on Drug Abuse. \textit{Drug use and addiction.} https://nida.nih.gov/
    \item Alcohol and Drug Foundation (ADF). \textit{Information about alcohol and other drugs.} https://adf.org.au/
    \item World Health Organization (WHO). \textit{Substance use and addiction.} https://www.who.int/
\end{itemize}

\clearpage
